\subsection{Probability Density Functions}
\hrulefill
\subsubsection*{Continuous Random Variables}
\paragraph{Definition:} A random variable $X$ is said to be \textbf{continuous} if its set of possible values is 
an interval or a collection of intervals on the real line. Examples of continuous random variables include:
\begin{itemize}
    \item The time required to complete a task.
    \item The height of a randomly selected person.
    \item The amount of rainfall in a given month.
    \item The weight of a randomly selected object.
\end{itemize}

\subsubsection*{Probability Density Function (PDF)}
\paragraph{Definition:} Let $X$ be a continuous random variable. Then a probability distrubution or probability 
density function (PDF) of $X$ is a function $f(x)$ such that for any two 
\[
    P(a \leq X \leq b) = \int_{a}^b f(x) dx
\]

Thus the probability that $X$ takes on a value between $a$ and $b$ is represented by the area under the curve of $f(x)$ from $a$ to $b$. 

\paragraph{Properties of PDF:}
\begin{itemize}
    \item $f(x) \geq 0$ for all $x$.
    \item $\int_{-\infty}^{\infty} f(x) dx = 1$.
    \item $P(X = x) = 0$ for any specific value of $x$.
    \item $P(X < x) = P(X \leq x)$.
\end{itemize}

\paragraph{Calculus Concepts Needed:}
\begin{itemize}
    \item Differentiation and Integration
    \item Polynomial and Exponential Functions
    \item Possibly Product and Quotient Rules
    \item Chain Rule
    \item Integration by Parts
    \item L'hopital's Rule
    \item $\int_0^\infty f(x) dx = \lim_{b \to \infty} \int_0^b f(x) dx$
\end{itemize}

\paragraph*{Example:}
Let $X$ be a continuous random variable with PDF:
\[
    f(x) = \begin{cases}
        cx & 0 \leq x \leq 3 \\
        0 & \text{otherwise}
    \end{cases}
\]
a. Find the value of $c$.\\
b. Compute $P(2 \leq X \leq 4)$\\

\paragraph*{a.}
\begin{align*}
    1 &= \int_{-\infty}^{\infty} f(x) dx \\
    1 &= \int_0^3 cx dx \\
    1 &= c \left[ \frac{x^2}{2} \right]_0^3 \\
    1 &= c \cdot \frac{9}{2} \\
    c &= \frac{2}{9}
\end{align*}

\paragraph*{b.}
\begin{align*}
    P(2 \leq X \leq 4) &= \int_2^4 f(x) dx \\
    &= \int_2^3 \frac{2}{9} x dx + \int_3^4 0 dx  & \text{(since $f(x) = 0$ for $x > 3$)}\\
    &= \frac{2}{9} \left[ \frac{x^2}{2} \right]_2^3 + 0 \\
    &= \frac{2}{9} \left( \frac{9}{2} - 2 \right) \\
    &= \frac{5}{9}
\end{align*}

\pagebreak

\subsubsection*{The Uniform Distribution}
\hrulefill

In the beginning of this course we considered smaple spaces in which every outcome is equally likely.
For example, let $X =$ the value of a fair six-sided die.
This is called the uniform distribution. The analogous continuous distribution is called the continuous 
uniform distribution. We have to consider all possible values in an interval $(a,b)$. The probability that $X$
falls into a subinterval should only depend on the length (and not the location) of the subinterval $(a,b)$.
This means that the probability density function must be constant. Let $f(x) = c \forall x \in (a,b)$.
Then we have
\[\int_a^b f(x) \, dx = 1 \implies c(b-a) = 1 \implies c = \frac{1}{b-a}.\]
Thus the probability density function of a continuous uniform distribution is
\[f(x) = \begin{cases}
\frac{1}{b-a} & a < x < b \\
0 & \text{otherwise}
\end{cases}.\]

\paragraph*{Notation:} $X \sim U(a,b)$ means that $X$ is a continuous uniform random variable on the 
interval $(a,b)$.

\paragraph*{CDF:} The CDF of $X \sim U(a,b)$ is
\[F(x) = P(X \leq x) = \begin{cases}
0 & x \leq a \\
\frac{x-a}{b-a} & a < x < b \\
1 & x \geq b
\end{cases}.\]

\paragraph*{Mean and Variance:} If $X \sim U(a,b)$, then
\[\mu = E(X) = \frac{a+b}{2}, \quad \sigma^2 = Var(X) = \frac{(b-a)^2}{12}.\]

\pagebreak

\paragraph*{Example:}
Suppose you arrie at a bus stop at 10:00 am. The bus will arrive at a time T uniformly distributed between 
10:00 am and 10:30 am.\\
a. What is the probability that you will have to wait more than 5 minutes?\\
b. If at 10:15 am the bus has not arrived, what is the probability that you will have to wait 
at least another 5 minutes?\\
c. What is the probability that the bus will arrive exactly at 10:15 am?

\paragraph*{a.}
\begin{align*}
    &X \sim U(0,30)& \\
    P(X>5) &= 1 - P(X \leq 5) = 1 - F(5) = 1 - \frac{5-0}{30-0} = \frac{25}{30} = \frac{5}{6}
\end{align*}

\paragraph*{b.}
\begin{align*}
    &X \sim U(0,30)& \\
    P(X>20 | X>15) &= \frac{P(X>20 \cap X>15)}{P(X>15)} = \frac{P(X>20)}{P(X>15)} = \frac{1 - F(20)}{1 - F(15)} \\
    &= \frac{1 - \frac{20-0}{30-0}}{1 - \frac{15-0}{30-0}} = \frac{\frac{10}{30}}{\frac{15}{30}} = \frac{2}{3}
\end{align*}

\paragraph*{c.}
\[P(X=15) = \int_{15}^{15} f(x) \,dx = 0.\]
This is because the probability of a continuous random variable taking on any single value is 0.
